\section{n2. Distinct MARKED normalization centres obstruct geometric envelopes, even for DOMAIN normalization.}
\textit{Supports: G331, n1.}

Marked-centre obstruction, proved relative to ENTIRE qualified G331 and n1. Let R be excellent reduced local of positive dimension, a a proper primary gauge, S its finite normalization, and rho the FULL normalized boundary quotient. Write M\_1,...,M\_q for S's finitely many closed centres over m and S\_i=S\_(M\_i), using the appropriate normal-domain factor of S. If one rho-fibre contains alpha on centre M\_1 and beta on a DISTINCT centre M\_2, there is a bounded strictly negative continuous signed PRIMARY-MODEL obstacle with neither a greatest geometric C nor a greatest geometric D minorant. S may be a DOMAIN; distinct minimal components are not required. No source CE or higher-dimensional resolution is assumed.

The local primary gauge a\_i=a S\_i has M\_i\textasciicircum{}r\_i contained in a\_i for some integer r\_i\textgreater{}=1. Every normalized w on that centre has w(M\_i)\textgreater{}=1/r\_i. Denominators outside M\_i are units and have value0, so these inequalities hold on ALL source supports, and w(a)=w(a\_i)=1. Chinese remainder gives h in S with h=1 modulo M\_1 and h=0 modulo EVERY other M\_i, and k in S with k=0 modulo M\_1 and k=1 modulo M\_2. At M\_1, h is a unit; elsewhere it is a nonunit. At M\_2, k is a unit. These are elements of the finite semilocal normalization, not falsely asserted downstairs units or idempotents.

Choose a regular c\_0 in the conductor c=(R:S). Since c is an S-ideal, c\_0 times EVERY element of S lies in R. A nontrivial fibre is conductor-infinite by G331, so alpha and beta kill c\_0. The conductor is proper and c\_0 is a nonunit of R; hence it is positive and nonzero on every generic centred source valuation. On M\_2 choose a normalized centred divisorial nu, available by G331 divisorial density. Its normal-domain order bound gives finite K\textgreater{}=1 such that nu(I)\textless{}=K beta(I) for EVERY unit-or-primary ideal of S\_2. More explicitly, after normalizing nu, nu(f)\textless{}=C\_nu ord\_(M\_2)(f) for nonzero nonunits and beta(M\_2)\textgreater{}=1/r\_2. Choose a generator attaining the primary ideal's finite beta-value. Then nu(I)\textless{}=C\_nu r\_2 beta(I); units are zero. Uniform comparison and actual decreasing infima consequently give g(rho beta)\textless{}=g(rho nu)/K for EVERY target geometric C or D function, including bottom. This bound is applied to the excellent NORMAL local domain S\_2; no analytic irreducibility of nonnormal R is asserted.

Put d=nu(c\_0)\textgreater{}0. Choose integers L\textgreater{}d and B\textgreater{}0 with D\_0=1+B(1-d/L)\textgreater{}K; write b=B+1. Choose an integer M\textgreater{}=L r\_1 and set ell=c\_0 k\textasciicircum{}M, an ACTUAL element of R. On EVERY point of centre M\_1, w(ell)\textgreater{}=M/r\_1\textgreater{}=L, while nu(ell)=d since k is a unit at M\_2. The ideal I=(ell)+a\textasciicircum{}L is proper primary: it contains a\textasciicircum{}L and ell cannot be a unit of R since its M\_1 value is positive. Define F=(B/L)v(I)-b v(a). On V\_a, F lies in [-b,-1], is exactly -1 on the ENTIRE M\_1 boundary image and at the common fibre, and F(rho nu)=-D\_0\textless{}-K. Primary evaluations prove continuity on every support. It is a signed exceptional REAL Cartier evaluation on a directed common ACTUAL primary blowup of I and a.

For each n\textgreater{}=1 choose an integer M\_n\textgreater{}=B n max\_(i!=1)r\_i. Put j\_n=c\_0 h\textasciicircum{}(M\_n) in R, T\_n=a\textasciicircum{}(bn)+(j\_n)a\textasciicircum{}n and q\_n=-v(T\_n)/n. T\_n is ACTUAL proper primary. On its actual blowup, -E\_(T\_n)/n is negative effective and relatively ample, so has positive degree on EVERY contracted proper integral curve BEFORE any limit. On M\_1 its exact profile is max(-b,-1-w(c\_0)/n), since h is a unit there. At every other centre, w(j\_n)\textgreater{}=M\_n/r\_i\textgreater{}=Bn, so q\_n=-b. These formulas include infinite c\_0-values and all nonminimal supports. Thus q\_n\textless{}=F EVERYWHERE. On the M\_1 finite-c\_0 locus, sup\_n q\_n=-1. That locus is dense by G331's all-support divisorial density: c\_0 is nonzero in S\_1, so every generic-support centred divisorial point has finite c\_0-value. In particular alpha is in its closure. A convergent NET suffices; no metrizability is used.

Every C/D minorant g\textless{}=F has g(rho beta)\textless{}=-D\_0/K\textless{}-1. If there were a greatest minorant, it would dominate all q\_n and equal -1 on the dense finite-c\_0 M\_1 locus, because F=-1 there. Its normalization pullback is USC. A strict sublevel neighbourhood of alpha at the forced value below -1 would meet that dense locus, a contradiction. Therefore neither greatest minorant exists. Equivalently the raw supremum p is -1 on that locus and \textless{}=-D\_0/K at the common point, so p is not USC. This proves the proposed CRT route, including actual ring membership, uniform centre bounds, proper primary models and every-support positivity.

There are quantitative consequences. Set eta=D\_0/K-1\textgreater{}0. Each canonical h\_m of n1 is continuous and \textless{}=p. At alpha it is \textless{}=-D\_0/K; its values on a finite-c\_0 M\_1 net approaching alpha tend to that value, whereas p on that net is -1. Therefore ||p-h\_m||\textgreater{}=eta in EVERY degree. The canonical algebra is not finitely generated by n1's proved implication. No exact formula for the whole raw p away from this locus and this fibre is claimed. Failure is stable for ||F'-F||\textless{}epsilon\textless{}(D\_0-K)/(K+1): q\_n-epsilon are C minorants of F', forcing a greatest minorant to be \textgreater{}=-1-epsilon on the dense locus, while the normal bound forces its common-fibre value \textless{}=(-D\_0+epsilon)/K\textless{}-1-epsilon. All functions are homogeneous on rays; primary changes divide by the same bounded positive denominator, preserve nonexistence and transform the canonical ideals exactly.

An algebraic sufficient condition is that c has a quotient-boundary fibre whose lifts lie in two distinct marked centres. In particular positive-dimensional intersections of distinct minimal components give such fibres: a valuation of the positive-dimensional quotient R/(p\_1+p\_2) lifts separately through the two finite normalizations. This recovers the previous crossing result as a special case rather than restating it as new mathematics. Universal CE implies that EVERY normalization fibre stays within ONE marked centre. This necessary condition does not assert that such a fibre must be singleton. Irreducible excellent domains whose normalization has several marked centres are covered whenever a full conductor fibre connects them; mere existence of several centres, without such a fibre, is insufficient. Curves remain consistent: their conductor quotient has empty centred boundary.

Same-marked-centre fibres, general NORMAL dimension\textgreater{}=3 CE, products, energy and full finite operators are not solved by this theorem. All model/source/audit/radial qualifications of G331 remain. Artifact n2.md contains the universal proof; n2.tex is its expanded delivery form. Rerun: python assemble.py; python verify\_delivery.py. Finite arithmetic is diagnostic, not certification of this theorem.

\[
 \ell=c_0k^M\in R,\qquad j_n=c_0h^{M_n}\in R,\qquad
 T_n=\mathfrak a^{(B+1)n}+(j_n)\mathfrak a^n.
\]
\[
 F=\frac BL v((\ell)+\mathfrak a^L)-(B+1)v(\mathfrak a),\quad
 g(\rho\beta)\le g(\rho\nu)/K<-1.
\]

